%Notes on section: I think I should x-ref to where I discuss Fowler in more depth. E.g. where I'm talking about the nature of objective demand. 

%\section{Fowler [better name needed!] \label{sec:fowler}}

Christopher Fowler (two papers here) has attempted to recreate the core model in agent form. It is useful to consider the approach he took and the obstacles discovered along the way. Fowler firstly attempted to take the core model out of its equilibrium assumptions; an agent approach, he argues, is a good method for achieving this. He set himself a specific goal: to ``re-create as exactly as possible the relationships in the economic [core] model'' (p.266) in agent form, meaning that ``to the extent possible, the equations used to express the analytic model have been maintained.'' (p.272) Fowler's direction of travel, then, is very much from the analytic core model towards agent modelling. For instance, he keeps the final set of simultaneous equations outlined in section \ref{sec:coremodelapx}. The main conclusion from his first attempt was that -

\begin{quote}
``although the equations of the analytic model can be mimicked in a way that is sufficient for a simulation to run, such a simulation cannot be made logically defensible without significantly altering the relationships among workers, firms and cities posited in the analytic model.'' (p.282)
\end{quote}

In initially binding himself to the assumptions of the analytic core model, Fowler was stuck with one of the more thorny abstractions: there are no firms explicitly defined. The equilibrium number of firms - a key feature of of the Dixit-Stiglitz \gls{monocomp} approach, as discussed in section \ref{sec:marketstructure} - is purely a consequence of consumer demand combined with love of variety. As he puts it, ``firms appear and disappear in cities based on the full employment of the workforce even to the point of following the lead of workers who move from place to place in order to benefit from increased real wages.'' (F2 p.2) 

\citealt{fowler_taking_2007} defined firm agents, but in sticking with the analytic model's mechanisms, avoided defining any mechanism for matching firms to workers. Following their own motivations, ``when firms move, they tend to want to move to places where wages are low; in contrast, workers tend to want to locate where wages are high'' (\citealt{fowler_taking_2007} p.276). 

Fowler's first pass at an agent version of the core model, then, lacking any method for linking firm and consumer behaviour even implicitly, could not produce equilibrium firm levels. Fowler notes that ``as a result, the model fails to move towards any of the equilibrium conditions predicted by its analytic counterpart.'' (p.266) In his second paper, Fowler attempts to address this by defining mechanisms, having already identified a list: 

\begin{quote}
``firms would bid for labor through wages and set prices for their goods at a level where they would expect to make the greatest profit. Workers would then consume based on the available prices and quantities; utilities and profits could be calculated for all agents. What sets this type of specification apart, however, is that it would make no assumptions about individual firms making the correct choices. Firms that chose correctly would be rewarded, and those that chose poorly could be dropped from the model.'' (\citealt{fowler_taking_2007} p.280-1)
\end{quote}
%The equations for things like production and wage would be truly independent rather than derivative so the model could handle a wide range of parameter configurations. 
%Actors with motivations, mechanisms connecting them and a variety of assumptions about the dynamics that these two together will produce. Fowler suggests some of the mechanisms that a full agent version of the core model might do: 

This is a fully agent-based picture he paints: each actor making decisions entirely independently, based on the information at their disposal. Fowler goes on to say, ``as complex as this specification sounds, it is actually relatively simple in an agent-based framework where bounded rationality, learning by doing and other types of decision-making all have substantial supporting bodies of literature.'' (ibid.) 

Not so simple. How to say? Break the problem into two components: first, how actors make decisions; second, whether they can collectively coordinate their actions into globally desirable economic outcomes. That, in a nutshell, is the `general' problem in general equilibrium analyses. 



Fowler discusses some of the dependencies that must be solved if an agent model is to work: ``the amount a firm can offer in wages is dependent on the amount it can produce and the price it receives for its goods. These quantities are affected, in turn, by consumer's wages, which depend on which firm employs them (and whether or not they have found employment at all.)'' (\citealt{fowler_taking_2007} p.277) 


Creating a functioning agent model of work and consumption can often involve making coding compromises. Fowler's approach to consumption reflects this: a tactical choice is made that avoids the need to consider rationing of limited goods. As Fowler puts it, ``workers select their optimal bundle of goods without reference to the actual supply of the good.'' (f2 p.10) Thus, they still have an `objective' demand level. He notes a key informational advantage: this sends ``clear signals to firms about the actual level of demand for their goods at the current price'' (ibid.) 
%I do the same currently: infinite stock. 

This reflects the prime difficulty of producing coordination in an agent model of this sort. A proponent of the analytic approach to the core model might argue that this demonstrates the power of the analytic method, but this misses the point Fowler's work has made. Just because dynamics of coordination are hidden behind equilibrium assumptions does not mean they have gone away. It is still very much an open question: how can actors, unable to coordinate through any central mechanism, produce stable spatial economic outcomes? Indeed, is there a model than can achieve this, or are there reasons to believe it might not be possible?
%Referring to the metis arg here. Analytics may describe and endpoint, but the metis argument rules out being able to produce a model of the process. Discuss!



Notes the 'absence of rational expectations' (p.275) - ``The firm and worker agents have perfect information in the sense that they can know the outcomes of their decisions up to a point, but the simulation environment is dynamic and the conditions that they observe at the time they decide where to move will not remain constant as other firms and workers make their own decisions.'' (p.275) Discussed much here in relation to what rationality means in an agent context. Do my blind agents, able only to know when a random grop improved their lot, count as rational?


Later: ``To create such a model, firms would have to have some mechanism with which they could predict an appropriate production level for themselves and estimate the levels chosen by their competition in an environment where sufficient labor supply is not guaranteed. The delicate balance among the equations of the analytic model does not allow for the error and uncertainty that are necessary parts of this sort of prediction, and so a new set of relationships needs to be specified.'' (p.282)
He notes that analytic choices made to gain tractability can ``actually make the problem less tractable'' in an agent setting. (p.283)



%%%%%%%%%%%%%%%
% Fowler 2
%%%%%%%%%%%%%%%
Fowler's goal is still to find optima in his model: he wants to ``explore the capacity of an economic system to identify a stable Nash Equilibrium without it being enforced by assumption.'' (F2 p.14)%Nash eq reminder: everyone's found a strategy that they wouldn't want to change if they can see what everyone else is doing. 

%Fowler's little summary of the core model: ``the effects of economies of scale, the role of preferences for variety in driving trade between regions, and the push/pull relationship between these two factors in shaping patterns of agglomeration and dispersal in economic activity based on transportation costs.'' (f2 p.2)

%``The research presented here maps out, as nearly as possible, the differences and similarities between the behavior of the agent based model and its deductive counterpart.'' f2p.3: gotta be good for me to bounce off. 

His idea of leakage: defined as the difference from the analytic optima. Which is interesting, for comparison, but my version of things is so far from the original, that sort of comparison isn't really feasible. It is also buys into the idea that the objective optima described by the core model do indeed exist. Or perhaps not: his goal is explicitly a comparison. If the goal instead is, starting with ABM, to ask, how to answer these questions?%Whoa, what an awful sentence!

``leakage is not necessarily a weakness of the agent-based model, but an explanation for how suboptimal outcomes may propagate in real systems.'' (f2 p.3) Sequence is important here, who gets to act first. Also, I haven't seen him really discuss that GE already deals with suboptimal outcomes: equilibria don't need to be optimal. But are they suboptimal for different reasons? And do those differing reasons have very different implications for what we should believe about the real world?

Leakage, he argues, comes from using `indivisible agents' (later he says firms can only hire whole people, rather than amounts of time) and differences between agents making simultaneous vs sequenced decisions. There's not much on that here, so plenty to play on. 

%To fulfill its basic objective the model is designed to ascertain when conditions are in equilibrium (when supply equals demand) and when that equilibrium is stable. To do this the model calculates the total supply and the total demand on each turn. When demand is within 1\% of supply, the model is considered to be in equilibrium and a counter starts to record the values for total supply and demand. If these values remain constant through 100 turns, then the model is considered to be in a stable equilibrium and halted. If no equilibrium is identified after 3000 model steps, then it is assumed that no stable equilibrium will be found for this model. (F2 p.14)

%``The complexity of the equations necessary to specify a relationship between the supply and demand for labor is a persuasive argument for the use if a simplifying assumption such as general equilibrium. However, this complexity points to the significance of the issues overlooked when full employment and optimal distribution of firms are assumed rather than aspired to.'' (F2 14-15)
%Yes, that's a point I've been trying to make: simplifying can hide complexity, but it's still there. However, cf. UNSW stand, when does it matter? It's not automatically true that complexity should be included.






%%%%%%%%%%%%%%%%%%%%%%%%%%%%%%%%%
% Equilibrium
%%%%%%%%%%%%%%%%%%%%%%%%%%%%%%%%%
%\subsection{Equilibrium \label{sec:equilibrium_n_agents}}

The problem of markets clearing in space... cf bullock, networks. Fowler: ``Equilibrium is the point in system space where supply equals demand, markets clear, efficiency reigns and overall welfare is maximized.'' (p.267) 

(Ladley bullock paper, unpublished AFAIK: The effects of market structure on a heterogeneous evolving population of traders. No other references: they just briefly mention trading in space can't fit into auction model.)

Hmm: this is an interesting effect of space, isn't it? If all actors are making market decisions on the head of a pin, an auction-style clearing market presents itself as the natural way to think about it. But as soon as actors have location, that becomes impossible. In a situation where production costs are the same for all actors, nearest actors will automatically take precedence.



Points out that equilibrium means being able to set some equations to zero, and set some rates of change equal to each other. [267]

Bring in physical example. Probably needs own section?

Nice quote to bounce off: ``As a rule, geographers have less faith in the importance of equilibrium values for understanding how an economy works.'' (p.271)
Someone who did think them important - as both Krugman and Isard argue, though Krugman not in any depth: we don't believe they exist in the real world, only that the real world tries to reach them. It never succeeds, but in understanding the equilibria, we can understand the processes of change that are going on in the real economy.
Gosh, that was a good summary, I'll have that! Note: Fowler goes on to say ``The economy as it has been represented in geography is typically far too chaotic to be considered at rest or to be operating at some market-clearing point where all labor is employed, all production is consumed and utility is maximized (Plummer et al., 1999).'' (p.271-2) Don't know who that ref is, but back to the point I was just making: no-one claims that the real economy will reach a rest point, only that the models help us understand the equilibria it moves towards. Think ball on rubber sheet in high wind.

%%%%%%%%%%%%%%%%%%%%%%%%%%%%%%%%%%%
% CUTTINZ
%%%%%%%%%%%%%%%%%%%%%%%%%%%%%%%%%%%
%The result was that ``workers locate without regard to labour demand in a particular locaiont

%One element in particular that he focuses on is the equilibrium number of firms - a key feature of of the Dixit-Stiglitz \gls{monocomp} approach. As discussed in section \ref{sec:marketstructure}, firms are derived from consumer demand. They have no independent, decision-making existence as such: ``the number of firms is simply a function of the labour required to produce a fixed optimal quantity of goods.'' (p.270) Given that agent modelling demands that firms exist as separate entities, this presents an immediate challenge for Fowler's attempted translation to an agent model. 
%``� the model assumes that the number of firms in a city is defined based on the amount of labour available in that city and the amount of labour necessary to produce the optimal amount of a good.'' (p.271)

%The direction taken by this project goes the other way: start with agents and work towards the concepts used in \gls{ge}.
%... This thesis comes at the problem from the opposite direction, in two respects. 

%He suggests this finding should lead to caution about making any policy claims based on the core model. I rather think it should instead lead to saying: agent models are different from their analytic counterparts, and those differences have to be respected. You cannot build a model of consumers and firms  - with explicit firms - without specifying how they interact. (You could perhaps model agent consumers and deduce firm numbers somehow, without explicitly creating firms as agents�) If the analytic argument about that interaction does not translate, then one that fits with an agent approach must be constructed.

%One reasonable alternative to a set of equilibrium solutions is a model that incorporates bounded rationality and the possibility of error. In this type of model, 
%Later: ``To create such a model, firms would have to have some mechanism with which they could predict an appropriate production level for themselves and estimate the levels chosen by their competition in an environment where sufficient labor supply is not guaranteed. 
%I should try and end up with a list of aspects of interest. So, these are things that ABM particularly draws attention to. To include:
%- how demand should work. Fowler goes for objective budget.
%- mechanisms for: linking production to labour; wages to employment. (So says Fowler.)
%- equilibrium level of firms. The connection between demand andntry/exit

%Hey! The seques nicely into a discussion on equilibrium, doesn't it? Aaawesome-ee-doo-dah!

%\begin{quote}``The logical step to improve the functionality of this model is to develop equations that account for the rationing of wages and production, based on worker employment (for wages) and availability (for production)... the model must be reformulated from its most basic assumptions before such rationing can be incorporated to move forward in this highly desirable direction.'' (\citealt{fowler_taking_2007} p.277)\end{quote}

%This is, of course, exactly the kind of argument Hayek was attempting to counter: success in providing a model rationing system would end up saying nothing about the fundamental problem of how disaggregated actors can coordinate production. There are benefits, however, to the `subjective' approach of testing possible goods bundles against utility, parallel to the benefits of blind search: no prior knowledge of the state of the market is required. Actors can act based on the informatin presented to them.

%Fowler quotes Amin and Thrift calling for a ``move away from the long shadow of economics'' \fixref{(Amin and Thrift 2000)}. Again, the direction of this thesis is actually more towards some aspects of traditional economics. 





